\documentclass[10pt]{article}
\usepackage[T1]{fontenc}
\usepackage[utf8]{inputenc}
\usepackage[french]{babel}
\usepackage{amsmath,amssymb,mathrsfs,mathpazo}
\usepackage{geometry}
\geometry{a4paper,margin=2cm}
\usepackage{graphicx,array,tabularx,tabularray,colortbl,multirow,multicol,booktabs,enumitem,pifont,ulem,fancybox,diagbox,cancel}
\usepackage{pstricks,pstricks-add,pst-plot,pst-tree,pst-node,pst-text,pst-func,pst-math,pst-bspline,pst-3dplot,pst-eucl,pst-solides3d}
\let\pscalculate\undefined
\usepackage{pst-bezier}
\usepackage{tikz,pgfplots,tkz-tab}
\usetikzlibrary{arrows,arrows.meta,patterns,calc,positioning,trees,decorations.markings,shapes.geometric,shapes.misc,shapes.arrows,intersections,decorations.pathmorphing,decorations.pathreplacing,backgrounds,shadows,babel}
\pgfplotsset{compat=1.18}
\usepgfplotslibrary{fillbetween}
\usepackage[np]{numprint}
\usepackage{listings,makecell,pas-tableur,twoopt}
\let\barre\undefined
\usepackage[pstricks]{tablvar}
\definecolor{darkgreen}{RGB}{0,100,0}
\providecommand{\red}{\color{red}}
\providecommand{\blue}{\color{blue}}
\providecommand{\green}{\color{green}}
\providecommand{\black}{\color{black}}
\providecommand{\white}{\color{white}}
\providecommand{\gray}{\color{gray}}
\providecommand{\cyan}{\color{cyan}}
\providecommand{\magenta}{\color{magenta}}
\providecommand{\yellow}{\color{yellow}}
\providecommand{\violet}{\color{violet}}
\providecommand{\orange}{\color{orange}}
\providecommand{\euro}{\texteuro}
\providecommand{\EUR}{\texteuro}
\providecommand{\ieme}{\textsuperscript{e}}
\providecommand{\degres}{\ensuremath{^\circ}}
\providecommand{\degree}{\ensuremath{^\circ}}
\providecommand{\vv}[1]{\overrightarrow{#1}}
\providecommand{\Coord}[1]{#1}
\providecommand{\decosix}{}
\providecommand{\starredbullet}{\ensuremath{\bullet}}
\pagestyle{empty}
\setlength{\parindent}{0pt}
\usepackage{adjustbox}\DeclareUnicodeCharacter{2236}{:}\DeclareUnicodeCharacter{2190}{\ensuremath{\leftarrow}}
\providecommand{\R}{\mathbb{R}}
\providecommand{\N}{\mathbb{N}}
\providecommand{\Z}{\mathbb{Z}}
\providecommand{\Q}{\mathbb{Q}}
\providecommand{\D}{\mathbb{D}}
\providecommand{\C}{\mathbb{C}}
\providecommand{\e}{{\mathrm{e}}}
\providecommand{\vect}[1]{\overrightarrow{#1}}
\providecommand{\vectt}[1]{\overrightarrow{\mathrm{#1}}}
\providecommand{\vvt}[1]{\overrightarrow{\mathrm{#1}}}
\providecommand{\np}[2][]{{#2~#1}}
\providecommand{\vv}[1]{{\overrightarrow{#1}}}
\renewcommand{\d}{{\mathrm{d}}}
\providecommand{\degres}{{^\circ}}
\providecommand{\degre}{{^\circ}}
\providecommand{\blue}{\color{blue}}
\providecommand{\red}{\color{red}}
\providecommand{\enskip}{\quad}
\providecommand{\llbracket}{\lbrack\!\lbrack}
\providecommand{\rrbracket}{\rbrack\!\rbrack}
\providecommand{\euro}{{\text{€}}}
\providecommand{\no}{{\text{n°}}}
\providecommand{\interval}[2]{{[#1;#2]}}
\providecommand{\intervFO}[2]{{[#1;#2[}}
\providecommand{\dotfill}{\ldots}

\providecommand{\ds}{\displaystyle}
\providecommand{\textding}[1]{#1}
\providecommand{\Oij}{\(\mathrm{O};\vec{\imath},\vec{\jmath}\)}
\providecommand{\Oijk}{\(\mathrm{O};\vec{\imath},\vec{\jmath},\vec{k}\)}
\providecommand{\Ouv}{\(\mathrm{O};\vec{u},\vec{v}\)}
\providecommand{\intervFF}[2]{\ensuremath{[#1;#2]}}
\renewcommand{\euro}{\ifmmode\text{\texteuro}\else\texteuro\fi}
\DeclareUnicodeCharacter{00B0}{\ifmmode^{\circ}\else\textdegree\fi}
\DeclareUnicodeCharacter{20AC}{\ifmmode\text{\texteuro}\else\texteuro\fi}
% Énoncés : Éducation nationale. Transcription et corrigés : APMEP, Denis Vergès et contributeurs. Découpage et classement : CoopMaths, Sébastien Lozano, Éric Elter et contributeurs. Adaptation de la présentation : Maths Explorer.
% CC BY-SA : https://coopmaths.fr/www/about/
% https://coopmaths.fr/alea/static/bac/2021/tex/bac_2021_06_asie_7_cor.tex
\begin{document}
\noindent\textit{APMEP / CoopMaths. CC BY-SA. Présentation adaptée par Maths Explorer.}\par\medskip
\textbf{Commun à tous les candidats}

\medskip






\begin{center}
\psset{unit=0.5cm,arrowsize=2pt 4}
\begin{pspicture}(6,10)
\pspolygon(0.5,0.5)(4,0)(4,7.2)(0.5,7.7)
\uput[d](0.5,0.5){B} \uput[d](4,0){C} \uput[ul](4,7.2){G} \uput[ul](0.5,7.7){F} 
\psline(4,0)(5,1.4)(5,8.6)(4,7.2)
\uput[r](5,1.4){D} \uput[ur](5,8.6){H} \uput[d](1.5,1.9){A} \uput[dr](4.5,0.7){K}\uput[l](1.5,5.5){I}
\uput[r](5,7.3){L}
\psline(5,8.6)(1.5,9.1)(0.5,7.7)
\uput[u](1.5,9.1){E}
\psline[linestyle=dashed]{->}(1.5,1.9)(1.5,5.5)
\psline[linestyle=dashed]{->}(1.5,1.9)(0.5,0.5)
\psline[linestyle=dashed]{->}(1.5,1.9)(5,1.4)
\psline[linestyle=dashed](4.5,0.7)(1.5,1.9)(5,7.3)
\psline[linestyle=dashed](1.5,5.5)(1.5,9.1)
\psline(4.5,0.7)(5,7.3)
\end{pspicture}
\end{center}

\bigskip







\begin{enumerate}
\item 
Avec C(1~;~1~;~0) et D(0~;~1~;~0), on obtient K$\left(\frac{1}{2}~;~1~;~0\right)$, donc $\vect{\text{AK}}\left(\frac{1}{2}~;~1~;~0\right)$ et on a $\vect{\text{AL}}\left(0~;~1~;~\dfrac{3}{2}\right)$.
\item  
	\begin{enumerate}
		\item 
\starredbullet~$\vect{n} \cdot \vect{\text{AK}} = 3 - 3 + 0 = 0$ ;

\starredbullet~$\vect{n} \cdot \vect{\text{AL}} = 0 - 3 + 3 = 0$ : le vecteur $\vect{n}$ est donc orthogonal à deux vecteurs non colinéaires du plan AKL, il est donc orthogonal à ce plan ; c'est donc un vecteur normal à ce plan.
		\item 
On a donc $M(x~;~y~;~z) \in (\text{AKL}) \iff \vect{\text{A}M} \cdot \vect{n} = 0 \iff 6x - 3y + 2z + d = 0$, \, avec $d \in \R$  et comme A appartient à ce plan on a : $ 0 + 0 + 0 + d = 0$.

Conclusion : $M(x~;~y~;~z) \in (\text{AKL}) \iff 6x - 3y + 2z = 0$.
		\item 
La droite $\Delta$ contient D et a pour vecteur directeur $\vect{n}$, donc :

$M(x~;~y~;~z) \in \Delta \iff \text{il existe }\, t \in \R, \, \vect{\text{D}M} = t\vect{n} \iff \left\{\begin{array}{l c r}
x		&=&6t\\
y - 1	&=&-3t\\
z		&=&2t
\end{array}\right.,\, t \in \R \iff \left\{\begin{array}{l c r}
x		&=&6t\\
y		&=&1-3t\\
z		&=&2t
\end{array}\right.,\, t \in \R$.
		\item 
		Le point N est donc le point commun au plan (AKL) et à la droite $\Delta$, donc ses coordonnées $(x~;~y~;~z)$ vérifient le système :
		
		$\left\{\begin{array}{l c r}
x		&=&6t\\
y		&=&1-3t\\
z		&=&2t\\
6x - 3y + 2z = 0
\end{array}\right.,\, t \in \R \Rightarrow 6 \times 6t + (-3)\times (1 - 3t) + 2\times 2t = 0 \iff$

$36t - 3 + 9t + 4t = 0 \iff 49t = 3 \iff t = \dfrac{3}{49}$ ; en remplaçant dans les trois premières équations du système, on obtient :

\renewcommand\arraystretch{1.8}
$\left\{\begin{array}{l c r}
x		&=&6\times \dfrac{3}{49}\\
y		&=&1-3\times\dfrac{3}{49}\\
z		&=&2\times \dfrac{3}{49}
\end{array}\right. \iff \left\{\begin{array}{l c r}
x		&=&\dfrac{18}{49}\\
y		&=&\dfrac{40}{49}\\
z		&=&\dfrac{6}{49}
\end{array}\right.$.
\renewcommand\arraystretch{1}
Conclusion : N$\left(\dfrac{18}{49}~;~\dfrac{40}{49}~;~\dfrac{6}{49} \right)$.
	\end{enumerate}







\item 
	\begin{enumerate}
		\item Le triangle ADK est rectangle en D ; on a par définition AD = 1 et DK $=  \frac{1}{2}$.

Donc $\mathcal{A}(\text{ADK}) = \dfrac{\text{AD} \times \text{DK}}{2} = \frac{1}{4}$.

D'autre part $\text{DL} = \frac{3}{2}$, donc 

$\mathcal{V}(\text{ADKL}) = \dfrac{\frac{1}{4} \times \frac{3}{2}}{3} = \dfrac{1}{8}$.
		\item 
On a $\vect{\text{DN}}\left(\frac{18}{49}~;~\frac{40}{49} - 1~;~\dfrac{6}{49} \right)$, soit  $\vect{\text{DN}}\left(\frac{18}{49}~;~\frac{- 9}{49}~;~\frac{6}{49} \right)$, donc :

$\text{DN}^2 = \left(\frac{18}{49}\right)^2 + \left(\frac{-9}{49}\right)^2  + \left(\frac{6}{49}\right)^2 = \frac{18^2 + 9^2 + 6^2}{49^2} = \frac{324 + 81 + 36}{49^2} = \frac{441}{49^2} = \frac{21^2}{49^2} = \left(\frac{21}{49}\right)^2$. 

Donc DN $ = \dfrac{21}{49} = \dfrac{3}{7}$.
		\item 
	En prenant comme base le triangle AKL, on a :
	
$\mathcal{V}(\text{ADKL}) = \dfrac{\mathcal{A}(\text{AKL}) \times \text{DN}}{3}$, soit $\dfrac{1}{8} = \dfrac{\mathcal{A}(\text{AKL}) \times \frac{3}{7}}{3}$, d'où 

$\mathcal{A}(\text{AKL}) = 7 \times \dfrac{1}{8} = \dfrac{7}{8}$~(u. a.).
	\end{enumerate}
\end{enumerate}

\bigskip

\end{document}
